\section{Week 12:}

\subsection{6.5 Irrotational / Potential Flow}

\begin{conceptbox}[Assumptions for Potential Flow]
We assume the flow is:
\begin{enumerate}
    \item \textbf{Incompressible} ($\rho = \text{const}$)
    \item \textbf{Inviscid} ($\mu = 0$)
    \item \textbf{Irrotational} ($\nabla \times \vec{v} = \vec{0}$)
\end{enumerate}
\end{conceptbox}

Since the flow is irrotational ($\nabla \times \vec{v} = \vec{0}$), vector calculus dictates that the velocity field must be the gradient of a scalar field. Therefore, there exists a scalar function $\Phi$ such that:
\begin{equation}
    \nabla \Phi = \vec{v}
\end{equation}
where $\Phi$ is defined as the \textbf{Velocity Potential}.

Using the continuity equation for an incompressible fluid:
\begin{equation}
    \nabla \cdot \vec{v} = 0 \quad (\text{since } \rho = \text{const})
\end{equation}

Substituting $\vec{v} = \nabla \Phi$ into the continuity equation yields the \textbf{Laplace equation} for the velocity potential:
\begin{eqbox}{Laplace Equation}
    \begin{equation}
        \nabla^2 \Phi = 0
    \end{equation}
\end{eqbox}

\subsubsection{Momentum Balance (Navier-Stokes for Inviscid, Incompressible Flow)}

Starting with the Euler equation (Navier-Stokes assuming inviscid $\mu=0$ and incompressible $\rho = \text{const}$):
\begin{equation}
    \rho \left[ \frac{\partial \vec{v}}{\partial t} + \vec{v} \cdot \nabla \vec{v} \right] = \rho \vec{g} - \nabla p
\end{equation}

We utilize the following vector identity for the convective acceleration term:
\begin{equation}
    \vec{v} \cdot \nabla \vec{v} = \frac{1}{2} \nabla (\vec{v} \cdot \vec{v}) - \vec{v} \times (\nabla \times \vec{v})
\end{equation}

Since the flow is \textit{irrotational}, the vorticity is zero ($\nabla \times \vec{v} = \vec{0}$). The convective acceleration simplifies to:
\begin{equation}
    \vec{v} \cdot \nabla \vec{v} = \frac{1}{2} \nabla \lVert \vec{v} \rVert^2 = \frac{1}{2} \nabla v^2
\end{equation}

Assuming a conservative body force field where $\vec{g} = -g\hat{k}$ (gravity acts downwards), we can write $\rho \vec{g} = -\nabla (\rho g z)$. Substituting this into the momentum balance:
\begin{equation}
    \rho \left[ \frac{\partial \vec{v}}{\partial t} + \frac{1}{2} \nabla v^2 \right] + \nabla p + \nabla (\rho g z) = \vec{0}
\end{equation}

Now, substituting $\vec{v} = \nabla \Phi$, we can rewrite the unsteady term:
\begin{align*}
    \rho \frac{\partial}{\partial t}(\nabla \Phi) + \frac{\rho}{2} \nabla v^2 + \nabla p + \nabla(\rho g z) &= \vec{0} \\
    \nabla \left[ \rho \frac{\partial \Phi}{\partial t} + \frac{\rho}{2} v^2 + p + \rho g z \right] &= \vec{0}
\end{align*}

Integrating over the spatial domain gives the \textbf{Unsteady Bernoulli Equation}:
\begin{eqbox}{Unsteady Bernoulli Equation}
    \begin{equation}
        \rho \frac{\partial \Phi}{\partial t} + \frac{\rho}{2} v^2 + p + \rho g z = \text{constant}
    \end{equation}
\end{eqbox}

If the flow is \textbf{steady}, then $\frac{\partial \Phi}{\partial t} = 0$, recovering the classic Bernoulli equation:
\begin{equation}
    \frac{\rho}{2} v^2 + p + \rho g z = \text{constant}
\end{equation}

\begin{mainbox}{Solution Strategy for Incompressible, Inviscid, Irrotational Flow}
\begin{itemize}
    \item \textbf{Step 1:} Solve the Laplace equation $\nabla^2 \Phi = 0$ with appropriate Boundary Conditions (BCs).
    \item \textbf{Step 2:} Compute the velocity field using $\vec{v} = \nabla \Phi$.
    \item \textbf{Step 3:} Use the Bernoulli equation to compute the pressure field $p$.
    \item \textbf{Step 4:} Integrate the pressure over boundaries to get the net force on a body.
\end{itemize}
\end{mainbox}

\newpage
\subsection{6.6 2D Flow - Stream Function}

From the conservation of mass ($\rho = \text{const}$), we have $\nabla \cdot \vec{v} = 0$.
In 2D Cartesian coordinates, this expands to:
\begin{equation}
    \frac{\partial u}{\partial x} + \frac{\partial v}{\partial y} = 0
\end{equation}

If a stream function $\Psi(x, y)$ exists with the properties:
\begin{equation}
    u = \frac{\partial \Psi}{\partial y}, \quad v = -\frac{\partial \Psi}{\partial x}
\end{equation}

Let's check continuity:
\begin{equation}
    \nabla \cdot \vec{v} = \frac{\partial u}{\partial x} + \frac{\partial v}{\partial y} = \frac{\partial}{\partial x}\left(\frac{\partial \Psi}{\partial y}\right) + \frac{\partial}{\partial y}\left(-\frac{\partial \Psi}{\partial x}\right) = \frac{\partial^2 \Psi}{\partial x \partial y} - \frac{\partial^2 \Psi}{\partial y \partial x} = 0
\end{equation}
$\implies$ A 2D velocity field derived from a stream function $\Psi$ will, by construction, automatically satisfy the continuity equation.

\subsubsection{Physical Meaning of the Stream Function}
Consider the exact differential of $\Psi$:
\begin{equation}
    d\Psi = \frac{\partial \Psi}{\partial x}dx + \frac{\partial \Psi}{\partial y}dy = -v\,dx + u\,dy
\end{equation}

For a constant value of $\Psi$ (i.e., along a line where $\Psi = \text{const}$), we have $d\Psi = 0$:
\begin{equation}
    -v\,dx + u\,dy = 0 \implies v\,dx = u\,dy \implies \frac{v}{u} = \frac{dy}{dx}
\end{equation}

This is exactly the equation of a \textbf{streamline} (where the velocity vector is tangent to the path). Thus, $\Psi$ is constant along a streamline.

\begin{notebox}[Polar Coordinates Extension]
For a 2D flow in polar coordinates $(r, \theta)$:
\begin{equation*}
    \vec{v} = v_r \hat{e}_r + v_\theta \hat{e}_\theta
\end{equation*}
The continuity equation is $\nabla \cdot \vec{v} = \frac{1}{r}\frac{\partial}{\partial r}(r v_r) + \frac{1}{r}\frac{\partial v_\theta}{\partial \theta} = 0$.
The stream function $\Psi(r, \theta)$ is defined as:
\begin{equation}
    v_r = \frac{1}{r}\frac{\partial \Psi}{\partial \theta}, \quad v_\theta = -\frac{\partial \Psi}{\partial r}
\end{equation}
\end{notebox}

\subsection{6.7 2D Potential Flow}

Assuming a 2D fluid flow that is \textbf{incompressible} and \textbf{inviscid}:
\begin{itemize}
    \item 2D flow $\implies$ There is a stream function $\Psi$.
    \item Irrotational $\implies$ There is a velocity potential $\Phi$.
\end{itemize}

\textbf{Q: What is the relation between the stream function $\Psi$ and velocity potential $\Phi$?}

For $\Phi$, mass conservation requires:
\begin{equation}
    0 = \nabla \cdot \vec{v} = \nabla \cdot (\nabla \Phi) = \nabla^2 \Phi \implies \boxed{\frac{\partial^2 \Phi}{\partial x^2} + \frac{\partial^2 \Phi}{\partial y^2} = 0}
\end{equation}

For $\Psi$, the irrotationality condition requires:
\begin{align}
    \vec{0} = \nabla \times \vec{v} &\implies 0 = \frac{\partial v}{\partial x} - \frac{\partial u}{\partial y} \\
    0 &= \frac{\partial}{\partial x}\left(-\frac{\partial \Psi}{\partial x}\right) - \frac{\partial}{\partial y}\left(\frac{\partial \Psi}{\partial y}\right) \\
    0 &= -\frac{\partial^2 \Psi}{\partial x^2} - \frac{\partial^2 \Psi}{\partial y^2} \implies \boxed{\frac{\partial^2 \Psi}{\partial x^2} + \frac{\partial^2 \Psi}{\partial y^2} = 0}
\end{align}

\begin{resultbox}
Both the Velocity Potential $\Phi$ and the Stream Function $\Psi$ satisfy the 2D Laplace equation.
\end{resultbox}
\newpage
\subsubsection{Orthogonality of $\Phi$ and $\Psi$}

Consider the differential of $\Phi$:
\begin{equation}
    d\Phi = \frac{\partial \Phi}{\partial x}dx + \frac{\partial \Phi}{\partial y}dy = u\,dx + v\,dy
\end{equation}

Along an \textbf{equipotential line} ($\Phi = \text{const}$), $d\Phi = 0$:
\begin{equation}
    u\,dx + v\,dy = 0 \implies \left.\frac{dy}{dx}\right|_{\Phi = \text{const}} = -\frac{u}{v}
\end{equation}

Recall that for a \textbf{streamline} ($\Psi = \text{const}$):
\begin{equation}
    \left.\frac{dy}{dx}\right|_{\Psi = \text{const}} = \frac{v}{u}
\end{equation}

Since the product of their slopes is $\left(-\frac{u}{v}\right) \times \left(\frac{v}{u}\right) = -1$, the equipotential lines ($\Phi = \text{const}$) are \textbf{perpendicular} to the streamlines ($\Psi = \text{const}$).

Furthermore, the closer the equipotential lines are, the higher the gradient $\nabla \Phi$, which implies the faster the fluid element is moving.
\vspace{16pt}

\begin{center}

\begin{tikzpicture}[scale=1.5]
    % Streamlines (Red)
    \foreach \y in {-1.5, -0.5, 0.5, 1.5} {
        \draw[accent, thick] (-3, \y) to[out=0, in=170] (1, \y-0.2) to[out=-10, in=160] (4, \y-0.8);
    }
    
    % Equipotential lines (Blue)
    \foreach \x in {-2, -0.5, 1, 2.5} {
        \draw[primary, thick] (\x, 2) to[out=-80, in=100] (\x+0.5, -2.5);
    }
    
    % Airfoil profile
    \filldraw[fill=white, draw=black, ultra thick] 
        (0.5,0.4) .. controls (0.5,1.0) and (2.5,0.8) .. (3.5,-0.6) 
        .. controls (2.0,-0.1) and (0.5,-0.1) .. (0.5,0.4);

    % Labels
    \node[accent, right] at (4, 0.7) {Streamlines ($\Psi = c$)};
    \node[primary, below] at (3, -2.5) {Equipotentials ($\Phi = c$)};
\end{tikzpicture}
\end{center}

\begin{exbox}[Example: Stagnation Point Flow]
Given the velocity potential:
\begin{equation*}
    \Phi = \frac{a}{2}(x^2 - y^2)
\end{equation*}
Compute the stream function $\Psi$ (if it exists).

\textbf{1. Check if incompressible ($\implies \nabla^2 \Phi = 0$):}
\begin{equation*}
    \frac{\partial^2 \Phi}{\partial x^2} + \frac{\partial^2 \Phi}{\partial y^2} = a + (-a) = 0
\end{equation*}
Since the Laplacian is zero, the flow is incompressible, and a stream function exists.

\textbf{2. Define $\Psi$ using the Cauchy-Riemann relations:}
First, derive the velocity field from $\Phi$:
\begin{equation*}
    u = \frac{\partial \Phi}{\partial x} = ax, \quad v = \frac{\partial \Phi}{\partial y} = -ay
\end{equation*}
Now, equate these to the stream function derivatives:
\begin{equation*}
    \begin{cases}
        u = \frac{\partial \Psi}{\partial y} = ax \\
        v = -\frac{\partial \Psi}{\partial x} = -ay
    \end{cases}
\end{equation*}

Integrate the first equation with respect to $y$:
\begin{equation*}
    \Psi = axy + f(x)
\end{equation*}
Differentiate this with respect to $x$ and plug into the second equation:
\begin{equation*}
    \frac{\partial \Psi}{\partial x} = ay + f'(x) = ay \implies f'(x) = 0 \implies f(x) = C
\end{equation*}
Thus, $\Psi = axy + C$. We can set $C = 0$ (free gauge choice).
\begin{equation*}
    \Psi = axy
\end{equation*}


\textbf{3. Plot the streamlines and equipotential lines:}
\begin{itemize}
    \item Streamlines: $\Psi = \text{const} \implies axy = c \implies y = \frac{c}{ax}$ (Hyperbolas in quadrants $1, 3$ or $2, 4$)
    \item Equipotentials: $\Phi = \text{const} \implies x^2 - y^2 = c'$ (Hyperbolas opening on axes)
\end{itemize}
\end{exbox}

\begin{center}
\begin{tikzpicture}[scale=2.5]
    % Axes
    \draw[->] (-3,0) -- (3,0) node[right] {$x$};
    \draw[->] (0,-3) -- (0,3) node[above] {$y$};
    
    % Asymptotes for equipotentials
    \draw[dashed, gray!60] (-2.5,-2.5) -- (2.5,2.5);
    \draw[dashed, gray!60] (-2.5,2.5) -- (2.5,-2.5);
    
    \begin{scope}
    \clip (-3,-3) rectangle (3,3);
    
    % Streamlines xy = c (White/cyan in notes, here we use secondary color)
    \foreach \c in {0.5, 1.2, 2.0} {
        \draw[secondary, thick] plot[domain=0.1:3, samples=50] (\x, {\c/\x});
        \draw[secondary, thick] plot[domain=-3:-0.1, samples=50] (\x, {\c/\x});
        \draw[secondary, thick] plot[domain=0.1:3, samples=50] (\x, {-\c/\x});
        \draw[secondary, thick] plot[domain=-3:-0.1, samples=50] (\x, {-\c/\x});
    }
    
    % Equipotential lines x^2 - y^2 = c (Blue in notes, here we use primary color)
    \foreach \c in {0.5, 1.5, 2.5} {
        % x^2 - y^2 = c => x = sqrt(y^2+c)
        \draw[primary, thick] plot[domain=-2.8:2.8, samples=50] ({sqrt(\x*\x+\c)}, \x);
        \draw[primary, thick] plot[domain=-2.8:2.8, samples=50] ({-sqrt(\x*\x+\c)}, \x);
        
        % y^2 - x^2 = c => y = sqrt(x^2+c) (represented as x^2 - y^2 = -c)
        \draw[primary, thick] plot[domain=-2.8:2.8, samples=50] (\x, {sqrt(\x*\x+\c)});
        \draw[primary, thick] plot[domain=-2.8:2.8, samples=50] (\x, {-sqrt(\x*\x+\c)});
    }
    \end{scope}

    \node[secondary, fill=white, inner sep=1pt] at (2, 1.5) {$\Psi = \text{const}$};
    \node[primary, fill=white, inner sep=1pt] at (2, 0.4) {$\Phi = \text{const}$};
\end{tikzpicture}
\end{center}


\begin{expertbox}[General Concept for Solving 2D Potential Flow]
The foundational principle of 2D potential flow dictates that both the velocity potential and the stream function obey Laplace's equation:
\begin{equation}
    \nabla^2 \Phi = 0 \quad \text{and} \quad \nabla^2 \Psi = 0
\end{equation}
Because the Laplace equation is \textbf{linear} and \textbf{homogeneous}, solutions can be superposed. Complex flow fields around arbitrary bodies can be systematically constructed by linearly combining elementary potential flows (such as uniform flows, sources, sinks, doublets, and point vortices).
\end{expertbox}
